\section{The sunflower bound}\label{sec:intro}

How many sets can a family contain before some \(r\) of them meet in
exactly the same core? Such a configuration is a \emph{sunflower}.
More precisely, distinct sets \(A_1,\ldots,A_r\) form an \(r\)-sunflower
if \(A_i\cap A_j=C\) for every \(i\ne j\). Removing the core \(C\)
leaves pairwise disjoint petals. The core can be empty, as it is
for any \(r\) disjoint sets.

Erd\H{o}s and Rado asked whether bounding the number of elements in
each set forces a sunflower once the family is exponentially large.
The ground set may be arbitrarily large. Write \(M_r(k)\) for the
largest size of a family of distinct \(k\)-element sets containing
no \(r\)-sunflower.

\begin{theorem}\label{thm:main}
For each integer \(r\ge3\) there is a constant \(B_r\) such that every
family of distinct sets of size at most \(k\) containing no
\(r\)-sunflower has at most \(B_r^k\) members, for every integer
\(k\ge1\). In particular \(M_r(k)\le B_r^k\). With logarithms to
base two, one may take
\[
 B_r=2^{C H^4 L^{5/2}},\qquad H=\log_2(2r),\quad L=\log_2(2H),
\]
where \(C\) is an absolute constant.
\end{theorem}

The classical argument gives \(M_r(k)\le k!(r-1)^k\)~\cite{ER}.
Take a maximal collection of disjoint members. It has at most \(r-1\)
sets, and its union, of size at most \(k(r-1)\), meets every member.
